\section{Introduction}
\label{sec:introduction}

\IEEEPARstart{I}{nter-Process} Communication (IPC) serves as the foundational communication fabric of modern operating systems, high-performance computing (HPC) software pipelines, database engines, and microkernel architectures~\cite{silberschatz2018operating,kerrisk2010linux}. Within shared-memory symmetric multiprocessing (SMP) and multi-core systems, intra-node communication speed frequently governs overall execution scaling. As governed by Amdahl's Law, as parallel worker counts expand across modern heterogeneous cores, any sequential communication bottleneck or kernel serialization penalty rapidly degrades parallel efficiency~\cite{herlihy2020art}.

The Linux operating system provides a diverse spectrum of IPC primitives established across decades of Unix standards and POSIX evolution~\cite{posix2018standard,stevens1999unixipc}. These mechanisms range from classical unidirectional byte-stream abstractions (Anonymous Pipes and Named FIFOs) and structured kernel-managed message queues (System V and POSIX Message Queues) to memory-mapped zero-copy regions (POSIX Shared Memory synchronized via semaphores) and bidirectional networking sockets (Unix Domain Stream Sockets)~\cite{kerrisk2010linux,stevens2004unix}. Despite this rich diversity, systems programmers and performance engineers face a pervasive challenge: \textit{selecting the optimal IPC mechanism for a given workload remains non-trivial, heavily dependent on message payload size, process topology, synchronization overhead, and underlying hypervisor virtualization}~\cite{danyliuk2026classical,gregg2020systems}.

\subsection{Background and Motivation}
In modern data-intensive applications, communicating processes exchange messages spanning multiple orders of magnitude, from 64-byte discrete control events to megabyte-scale tensor buffers. Conventionally, applications adopt a single static IPC mechanism throughout their execution lifecycle. However, static binding forces severe performance compromises:
\textbf{1) Kernel Double-Copy Overhead:} Byte-stream mechanisms (Pipes, FIFOs, Sockets) and message queues copy data twice: from userspace to kernel memory, and from kernel memory to the receiver, causing memory bandwidth saturation and cache pollution for large payloads~\cite{drepper2007every}.
\textbf{2) Synchronization in Shared Memory:} Although POSIX Shared Memory achieves zero kernel data copies by mapping physical pages directly into communicating address spaces, it requires explicit userspace synchronization (POSIX semaphores or mutexes backed by Linux \texttt{futex})~\cite{franke2002futex}. Under scheduling contention, lock acquisition latency can offset memory copy advantages for small messages.
\textbf{3) Buffer Saturation and Flow Control:} Traditional pipes employ a default circular buffer capacity of 64~KB (16 pages on x86\_64 Linux)~\cite{bovet2005understanding}. Payloads exceeding this threshold saturate the buffer, forcing repeated blocking context switches that throttle throughput.

\subsection{Problem Statement and Research Gap}
While prior empirical studies, such as Danyliuk's investigation of classical Linux IPC mechanisms for HPC scenarios~\cite{danyliuk2026classical}, have established valuable baseline latency comparisons for small messages (64--1,024 bytes) on physical hardware, critical empirical dimensions remain unexplored. Specifically, existing investigations often defer large-payload throughput evaluation ($\ge 64$~KB), do not analyze tail-latency percentiles ($P_{95}, P_{99}$) or system-call overheads under virtualization, leave multi-process MPMC concurrency scaling unaddressed, and omit adaptive runtime transport selection. Furthermore, while custom kernel mailboxes offer theoretical low-overhead communication, their uncontended loopback latency floor relative to the operating system's multi-process context-switch floor ($95\text{--}130\,\mu\text{s}$ under virtualization) has not been empirically contrasted within a unified experimental framework.

\subsection{Research Objectives}
To address these technical gaps, this paper investigates five central research questions:
\textbf{RQ1 (Payload Sensitivity):} How does payload size (64~B to 256~KB) influence latency, throughput, and system-call frequency?
\textbf{RQ2 (Memory vs. Stream Trade-offs):} What are the quantitative trade-offs between zero-copy shared memory and kernel-buffered streaming across message scales?
\textbf{RQ3 (Adaptive Selection):} Can a lightweight, telemetry-driven selector dynamically route payloads to minimize latency regret?
\textbf{RQ4 (Kernel Mailbox Feasibility):} What performance benefits and limits emerge from a custom driver-assisted character-device mailbox?
\textbf{RQ5 (Concurrency Scalability):} How do classical IPC primitives behave when scaled from 1P:1C to contended 4P:4C topologies?

\subsection{Major Contributions}
This research introduces \textbf{Mini-A-IPC}, an open-source systems evaluation framework and adaptive communication architecture. The primary contributions of this paper are:
\textbf{1) Comprehensive Implementation Suite:} Complete, POSIX-compliant, secure (\texttt{0600} permissions), and race-free C11 implementations of six classical Linux IPC mechanisms: Anonymous Pipes, Named FIFOs, System V Message Queues, POSIX Message Queues, POSIX Shared Memory, and Unix Domain Stream Sockets.
\textbf{2) Adaptive Telemetry Controller:} Design and implementation of an adaptive transport selector combining static thresholding with dynamic Exponentially Weighted Moving Average (EWMA) latency telemetry and an 8\% hysteresis switching margin (\texttt{src/selector.c}), verified through automated unit tests, with empirical evaluation of the deterministic static threshold path demonstrating 0.66~ns decision overhead.
\textbf{3) Custom Linux Kernel Mailbox:} An in-kernel circular-buffer character device (\texttt{/dev/aipc}) featuring sleepable mutex synchronization, blocking wait queues, and an atomic reference-counting teardown barrier designed to mitigate module-unload Use-After-Free (UAF) race conditions.
\textbf{4) Exhaustive Empirical Evaluation:} Over 3.4 million controlled, randomized, and cache-primed round-trip message exchanges (recording 2.73 million valid timed exchanges and exceeding 4.5 million discrete message transfers) on primary Fedora 45 Linux (Kernel 7.3.3, KVM virtualization), collecting tail percentiles ($P_{50}, P_{95}, P_{99}$), classified syscall profiles, and OS resource telemetry.

\subsection{Organization of the Paper}
Section~\ref{sec:related_work} reviews related work. Section~\ref{sec:architecture} details system architecture. Section~\ref{sec:ipc_implementation} examines low-level IPC mechanics. Section~\ref{sec:adaptive_kernel} presents the adaptive telemetry selector and kernel driver. Section~\ref{sec:methodology} describes our experimental methodology. Section~\ref{sec:results} provides empirical performance results. Section~\ref{sec:discussion} analyzes concurrency, evaluates adaptive regret, and answers RQ1--RQ5. Finally, Section~\ref{sec:conclusion} concludes the paper.
